\documentclass[11pt]{beamer}

\title[]{Information Flow and the Role of Geometry in Graph Neural Networks: A Preliminary Study}
\author{Luigi Fracassetti}
\date{\today}
\usetheme{Madrid}
\usecolortheme{whale}
\setbeamertemplate{navigation symbols}{}
\setbeamertemplate{headline}{}
\usepackage{pgfplots}
\usepackage{graphicx}
\usepackage[english]{babel}
\usepackage{mathtools}
\usepackage{amssymb}
\usepackage{booktabs}
\setbeamercovered{transparent}
\usepackage{stackrel}
\pgfplotsset{compat=1.18}

\begin{document}

\begin{frame}
  \titlepage
  \vspace{-0.5cm}
  \begin{block}{Current status}
    The codebase now supports reproducible Neural Sheaf Diffusion experiments,
    baseline controls, and post-training diagnostics for long-range information
    flow on synthetic, transfer, and city-scale graphs.
  \end{block}
\end{frame}

\begin{frame}{Scientific Question and Hypothesis}
  \begin{block}{Question}
    Does learned sheaf geometry change how information travels across graph
    bottlenecks, compared with scalar or fixed-message-passing baselines?
  \end{block}

  \begin{itemize}
    \item Over-squashing is studied through source-to-target sensitivity:
    \[
      J_{v,u}=\frac{\partial z_v}{\partial x_u}.
    \]
    \item Geometry enters through learned restriction and transport maps on edges.
    \item Main comparison: \emph{performance} versus \emph{measured information flow}.
  \end{itemize}

  \begin{block}{Working hypothesis}
    Non-scalar, learned sheaves should preserve useful long-range signals by
    shaping anisotropic transport across critical regions of the graph.
  \end{block}
\end{frame}

\begin{frame}{Experimental Coverage Implemented}
  \begin{table}
    \centering
    \small
    \begin{tabular}{llll}
      \toprule
      Family & Dataset settings & Task & Depth \\
      \midrule
      Barbell & 1 setting & opposite-clique regression & 3 layers \\
      Transfer & 22 settings & source classification & shortest path \\
      Cities & Paris, Shanghai & eccentricity classes & 16 layers \\
      \bottomrule
    \end{tabular}
  \end{table}

  \begin{itemize}
    \item Benchmark grid: general, orthogonal, diagonal, and identity NSD.
    \item Stalk dimensions \(d\in\{2,3,5\}\), hidden dimensions \(h\in\{16,32\}\).
    \item Three benchmark seeds, giving \(24\) architectures and \(72\) runs per dataset setting.
    \item Analysis profile adds GCN, GAT, GraphSAGE, MLP, frozen orthogonal sheaves,
    scalar-stalk controls, and constant-total-width controls.
  \end{itemize}
\end{frame}

\begin{frame}{Diagnostics Now Available}
  \begin{columns}[T,totalwidth=\textwidth]
    \begin{column}{0.49\textwidth}
      \begin{block}{Influence}
        \[
          \|J_{v,u}\|_1,\quad
          \|J_{v,u}\|_F,\quad
          \sigma_{\max}(J_{v,u})
        \]
        Aggregated by graph distance using shell totals, shell means, bootstrap
        confidence intervals, decay slopes, and influence radius.
      \end{block}
    \end{column}
    \begin{column}{0.49\textwidth}
      \begin{block}{Pathwise Flow}
        Local layer Jacobians are multiplied along monotone geodesic walks:
        \[
          A_k(a,b)=
          \frac{\partial h_a^{(k)}}{\partial h_b^{(k-1)}}.
        \]
        This separates full, geodesic, canonical-path, residual, and cancellation terms.
      \end{block}
    \end{column}
  \end{columns}

  \vspace{0.2cm}
  \begin{block}{Geometry}
    Each layer stores restriction-map spectra, transport norms, condition numbers,
    effective ranks, bottleneck strengths, and learned effective Ollivier--Ricci curvature.
  \end{block}
\end{frame}

\begin{frame}{Engineering State and Next Steps}
  \begin{itemize}
    \item Unified CLI, configuration system, deterministic run IDs, SQLite index,
    checkpoint storage, summaries, and export commands are in place.
    \item GPU launchers exist for barbell, transfer, cities, and analysis; RunPod
    deployment is organized into barbell, three transfer shards, and one worker per city.
    \item Local smoke artifacts exist for barbell and transfer, including completed
    analysis summaries, geometry tables, influence tables, pathwise tables, and figures.
    \item City benchmark configuration is ready: Paris and Shanghai, \(16\) NSD layers,
    AdamW, \(2000\) epochs, patience \(200\), seeds \(0,1,2\).
  \end{itemize}

  \begin{alertblock}{Interpretation}
    The current repository state demonstrates that the pipeline runs end-to-end.
    Full scientific conclusions still require the complete benchmark and analysis runs.
  \end{alertblock}

  \begin{block}{Next}
    Run the complete benchmark grid, analyze initial and best checkpoints, and
    connect task metrics to influence decay, pathwise cancellation, and learned
    curvature changes.
  \end{block}
\end{frame}

\end{document}
